% GENERATED FILE. Edit canonical lessons, then run npm run generate:books.
% canonical-source-hash: fnv1a64:b77573ca1253ee04
% canonical-lessons: JA-C234-shiai, JA-C234-katsu, JA-C234-makeru, JA-C234-senshu, JA-C234-ouen

\chapter{Winning and Losing}
\label{ch:ja-winning-and-losing}
\hlchaptermodality{234}

\begin{chapteropening}
\textbf{By the end of this chapter:} \emph{I can say who won and who lost a match.}

Complete the last lesson of chapter 234: I can say who won and who lost a match.
\end{chapteropening}

\section[\texorpdfstring{\ja{しあい} (shiai)}{しあい (shiai)}]{\ja{しあい} (shiai) — a match; a game}

\label{lesson:JA-C234-shiai}



\begin{quote}
Before the new one: say the Japanese for to run, then the Japanese for exercise.
\end{quote}

\subsection*{You'll want to know: \ja{しあい}}
\textbf{\ja{しあい}} — \emph{shiai} — \textquotedblleft{}a match; a game\textquotedblright{}.

\begin{quote}
\textbf{\ja{せんしゅう} \ja{がっこうで} \ja{しあいが} \ja{ありました。}} — \emph{senshū gakkō de shiai ga arimashita} — \textquotedblleft{}There was a match at school last week.\textquotedblright{}
\end{quote}

\begin{cousinweb}
\textbf{\ja{しあい}} (\emph{shiai}) is built on native \emph{ai}, \textquotedblleft{}meeting, matching\textquotedblright{}, from \emph{au}, inherited rather than borrowed. JMdict gives two spellings whose first characters, read \emph{shi}, mean \textquotedblleft{}try\textquotedblright{} and \textquotedblleft{}serve\textquotedblright{}; which of them, if either, is the source of \emph{shi} is not settled (JMdict, KANJIDIC2).
\end{cousinweb}

\begin{grammarlens}[title={What you've built}]
One more word for everyday life at A2.
\end{grammarlens}

\subsection*{Your turn}
\begin{itemize}
\raggedright
  \item \emph{Say it:} \emph{shiai}
  \item \emph{Say it:} \emph{shiai}, once more
  \item \emph{Say it:} \emph{senshū gakkō de shiai ga arimashita} — There was a match at school last week.
  \item \emph{Recall:} say the Japanese for to run, then the Japanese for exercise, then say \emph{shiai} again
\end{itemize}

\begin{tcolorbox}[breakable,colback=teal!4,colframe=teal!35!black,title={Before you move on}]
What does \ja{しあい} mean? (\textquotedblleft{}A match; a game\textquotedblright{}.) Say it once more.
\end{tcolorbox}
\section[\texorpdfstring{\ja{かつ} (katsu)}{かつ (katsu)}]{\ja{かつ} (katsu) — to win}

\label{lesson:JA-C234-katsu}



\begin{quote}
Before the new one: say the Japanese for exercise, then the Japanese for a match.
\end{quote}

\subsection*{You'll want to know: \ja{かつ}}
\textbf{\ja{かつ}} — \emph{katsu} — \textquotedblleft{}to win\textquotedblright{}.

\begin{quote}
\textbf{\ja{きのうの} \ja{しあいで} \ja{あにが} \ja{かちました。}} — \emph{kinō no shiai de ani ga kachimashita} — \textquotedblleft{}My elder brother won yesterday's match.\textquotedblright{}
\end{quote}

Polite \textbf{\ja{かちます}}, \emph{kachimasu}; a godan verb.

\begin{cousinweb}
\textbf{\ja{かつ}} (\emph{katsu}) is a native Japanese word (wago), inherited rather than borrowed from Chinese; its kanji was attached later as a kun'yomi (JMdict, KANJIDIC2).
\end{cousinweb}

\begin{grammarlens}[title={What you've built}]
One more word for everyday life at A2.
\end{grammarlens}

\subsection*{Your turn}
\begin{itemize}
\raggedright
  \item \emph{Say it:} \emph{katsu}
  \item \emph{Say it:} \emph{katsu}, once more
  \item \emph{Say it:} \emph{kinō no shiai de ani ga kachimashita} — My elder brother won yesterday's match.
  \item \emph{Recall:} say the Japanese for exercise, then the Japanese for a match, then say \emph{katsu} again
\end{itemize}

\begin{tcolorbox}[breakable,colback=teal!4,colframe=teal!35!black,title={Before you move on}]
What does \ja{かつ} mean? (\textquotedblleft{}To win\textquotedblright{}.) Say it once more.
\end{tcolorbox}
\section[\texorpdfstring{\ja{まける} (makeru)}{まける (makeru)}]{\ja{まける} (makeru) — to be beaten; to lose (a game)}

\label{lesson:JA-C234-makeru}



\begin{quote}
Before the new one: say the Japanese for a match, then the Japanese for to win.
\end{quote}

\subsection*{You'll want to know: \ja{まける}}
\textbf{\ja{まける}} — \emph{makeru} — \textquotedblleft{}to be beaten; to lose (a game)\textquotedblright{}.

\begin{quote}
\textbf{\ja{せんしゅうの} \ja{しあいで} \ja{まけました。}} — \emph{senshū no shiai de makemashita} — \textquotedblleft{}We lost last week's match.\textquotedblright{}
\end{quote}

Polite \textbf{\ja{まけます}}, \emph{makemasu}; an ichidan verb.

\begin{cousinweb}
\textbf{\ja{まける}} (\emph{makeru}) is a native Japanese word (wago), inherited rather than borrowed from Chinese; its kanji was attached later as a kun'yomi (JMdict, KANJIDIC2).
\end{cousinweb}

\begin{grammarlens}[title={What you've built}]
One more word for everyday life at A2.
\end{grammarlens}

\subsection*{Your turn}
\begin{itemize}
\raggedright
  \item \emph{Say it:} \emph{makeru}
  \item \emph{Say it:} \emph{makeru}, once more
  \item \emph{Say it:} \emph{senshū no shiai de makemashita} — We lost last week's match.
  \item \emph{Recall:} say the Japanese for a match, then the Japanese for to win, then say \emph{makeru} again
\end{itemize}

\begin{tcolorbox}[breakable,colback=teal!4,colframe=teal!35!black,title={Before you move on}]
What does \ja{まける} mean? (\textquotedblleft{}To be beaten; to lose (a game)\textquotedblright{}.) Say it once more.
\end{tcolorbox}
\section[\texorpdfstring{\ja{せんしゅ} (senshu)}{せんしゅ (senshu)}]{\ja{せんしゅ} (senshu) — a player; an athlete}

\label{lesson:JA-C234-senshu}



\begin{quote}
Before the new one: say the Japanese for to win, then the Japanese for to be beaten.
\end{quote}

\subsection*{You'll want to know: \ja{せんしゅ}}
\textbf{\ja{せんしゅ}} — \emph{senshu} — \textquotedblleft{}a player; an athlete\textquotedblright{}.

\begin{quote}
\textbf{\ja{きょねん} \ja{あには} \ja{がっこうの} \ja{せんしゅでした。}} — \emph{kyonen ani wa gakkō no senshu deshita} — \textquotedblleft{}Last year my elder brother was one of the school's players.\textquotedblright{}
\end{quote}

\begin{cousinweb}
\textbf{\ja{せんしゅ}} (\emph{senshu}) is Sino-Japanese (kango): two Chinese characters read with their on'yomi, readings borrowed from Chinese. They are \emph{sen}, \textquotedblleft{}choose\textquotedblright{} (Middle Chinese \emph{sjwenX}) and \emph{shu}, \textquotedblleft{}hand\textquotedblright{} (\emph{syuwX}) (KANJIDIC2; Middle Chinese in Baxter's transcription).
\end{cousinweb}

\begin{grammarlens}[title={What you've built}]
One more word for everyday life at A2.
\end{grammarlens}

\subsection*{Your turn}
\begin{itemize}
\raggedright
  \item \emph{Say it:} \emph{senshu}
  \item \emph{Say it:} \emph{senshu}, once more
  \item \emph{Say it:} \emph{kyonen ani wa gakkō no senshu deshita} — Last year my elder brother was one of the school's players.
  \item \emph{Recall:} say the Japanese for to win, then the Japanese for to be beaten, then say \emph{senshu} again
\end{itemize}

\begin{tcolorbox}[breakable,colback=teal!4,colframe=teal!35!black,title={Before you move on}]
What does \ja{せんしゅ} mean? (\textquotedblleft{}A player; an athlete\textquotedblright{}.) Say it once more.
\end{tcolorbox}
\section[\texorpdfstring{\ja{おうえん} (ōen)}{おうえん (ōen)}]{\ja{おうえん} (ōen) — cheering; support}

\label{lesson:JA-C234-ouen}



\begin{quote}
Before the new one: say the Japanese for to be beaten, then the Japanese for a player.
\end{quote}

\subsection*{You'll want to know: \ja{おうえん}}
\textbf{\ja{おうえん}} — \emph{ōen} — \textquotedblleft{}cheering; support\textquotedblright{}.

\begin{quote}
\textbf{\ja{きのう} \ja{みんなで} \ja{せんしゅの} \ja{おうえんを} \ja{しました。}} — \emph{kinō minna de senshu no ōen o shimashita} — \textquotedblleft{}Yesterday we all cheered for the players.\textquotedblright{}
\end{quote}

\begin{cousinweb}
\textbf{\ja{おうえん}} (\emph{ōen}) is Sino-Japanese (kango): two Chinese characters read with their on'yomi, readings borrowed from Chinese. They are \emph{ō}, \textquotedblleft{}respond\textquotedblright{} (Middle Chinese \emph{'ingH}) and \emph{en}, \textquotedblleft{}help\textquotedblright{} (\emph{hjwenH}) (KANJIDIC2; Middle Chinese in Baxter's transcription).
\end{cousinweb}

\begin{grammarlens}[title={What you've built}]
One more word for everyday life at A2.
\end{grammarlens}

\subsection*{Your turn}
\begin{itemize}
\raggedright
  \item \emph{Say it:} \emph{ōen}
  \item \emph{Say it:} \emph{ōen}, once more
  \item \emph{Say it:} \emph{kinō minna de senshu no ōen o shimashita} — Yesterday we all cheered for the players.
  \item \emph{Recall:} say the Japanese for to be beaten, then the Japanese for a player, then say \emph{ōen} again
\end{itemize}

\begin{tcolorbox}[breakable,colback=teal!4,colframe=teal!35!black,title={Before you move on}]
What does \ja{おうえん} mean? (\textquotedblleft{}Cheering; support\textquotedblright{}.) Say it once more.
\end{tcolorbox}
